% Part: first-order-logic
% Chapter: proof-systems
% Section: natural-deduction

\documentclass[../../../include/open-logic-section]{subfiles}

\begin{document}

\iftag{FOL}
      {\olfileid[fr]{fol}{prf}{ntd}}
      {\olfileid[fr]{pl}{prf}{ntd}}

\olsection{La déduction naturelle}

La déduction naturelle est un système de !!{derivation} qui vise à
refléter le raisonnement effectif, en particulier le raisonnement
discipliné des mathématiciens. Celui-ci suit divers schémas « naturels ».
Par exemple, le raisonnement par cas permet d'établir une conclusion à
partir d'une prémisse disjonctive en montrant que cette conclusion se
déduit de chacun des deux termes de la disjonction. Le raisonnement par
l'absurde permet d'établir une conclusion en montrant que sa négation
mène à une contradiction. Pour démontrer une implication « si \dots
alors \dots », on montre que le conséquent se déduit de l'antécédent.
La déduction naturelle formalise certaines de ces inférences naturelles.
À chaque connecteur logique et à chaque quantificateur sont associées
deux règles, l'une d'introduction et l'autre d'élimination, qui
correspondent chacune à un tel schéma d'inférence. Ainsi,
$\Intro{\lif}$ correspond à la démonstration d'une implication, et
$\Elim{\lor}$ au raisonnement par cas. La règle $\Elim{\land}$ est
particulièrement simple : elle permet de passer de $!A \land !B$ à~$!A$
(ou à~$!B$).

La déduction naturelle se distingue notamment des autres systèmes de
!!{derivation} par son usage des hypothèses. Une !!{derivation} en
déduction naturelle est un arbre de !!{formula}s. Une seule !!{formula}
figure à la racine de cet arbre ; ses « feuilles » sont des !!{formula}s
à partir desquelles la conclusion est dérivée. En déduction naturelle,
certaines !!{formula}s situées aux feuilles jouent un rôle au cours de la
!!{derivation}, mais ne sont plus requises lorsque l'on atteint la
conclusion. Cela correspond à la pratique qui consiste, dans un
raisonnement effectif, à introduire des hypothèses qui ne restent en
vigueur que temporairement. Par exemple, dans un raisonnement par cas,
on suppose successivement vrai chacun des termes de la disjonction ;
pour démontrer une implication, on suppose vrai l'antécédent ; dans un
raisonnement par l'absurde, on suppose vraie la négation de la conclusion.
Introduire ainsi des hypothèses temporaires, puis s'en défaire pour
établir une étape intermédiaire, est caractéristique de la déduction
naturelle. Les formules situées aux feuilles d'une !!{derivation} en
déduction naturelle s'appellent des hypothèses, et certaines règles
d'inférence permettent de les « !!{discharge} ». Par exemple, si une
!!{derivation} établit~$!B$ à partir d'hypothèses parmi lesquelles
figure~$!A$, la règle $\Intro{\lif}$ permet d'inférer~$!A \lif !B$ et
de décharger toute hypothèse choisie de la forme~$!A$. Pour indiquer quelle
inférence décharge quelles hypothèses, on marque cette inférence et
les hypothèses qu'elle décharge d'un même numéro. Si les hypothèses qui
restent !!{undischarged}s à la fin de la !!{derivation} sont vraies, la
conclusion l'est également. Une !!{derivation} établit donc que sa
conclusion est une conséquence logique de ses hypothèses
!!{undischarged}s.

La relation $\Gamma \Proves !A$ fondée sur la déduction naturelle est
vérifiée si et seulement s'il existe une !!{derivation} dont le dernier
!!{sentence} est~$!A$ et dont toute feuille !!{undischarged} appartient
à~$\Gamma$. $!A$ est un théorème en déduction naturelle si et seulement
s'il existe une !!{derivation} dont le dernier !!{sentence} est~$!A$
et dont toutes les hypothèses sont !!{discharged}s. Par exemple, la
!!{derivation} suivante établit que $\Proves (!A \land !B) \lif !A$ :
\begin{prooftree}
  \AxiomC{$\Discharge{!A \land !B}{1}$}
  \RightLabel{\Elim{\land}}
  \UnaryInfC{$!A$}
  \DischargeRule{\Intro{\lif}}{1}
  \UnaryInfC{$(!A \land !B) \lif !A$}
\end{prooftree}
Le numéro~$1$ indique que l'hypothèse $!A \land !B$ est
!!{discharged} lors de l'inférence \Intro{\lif}.

Un ensemble~$\Gamma$ est incohérent si et seulement si
$\Gamma \Proves \lfalse$ en déduction naturelle. La règle \FalseInt{}
permet de !!{derive} n'importe quel !!{sentence} d'un ensemble incohérent.

Les systèmes de déduction naturelle ont été élaborés par Gerhard Gentzen
et Stanis\l{}aw Ja\'skowski dans les années 1930, puis développés
par Dag Prawitz et Frederic Fitch. Comme leurs inférences reflètent
des méthodes naturelles de démonstration, ils ont la faveur des
philosophes. Les versions de Fitch sont souvent utilisées dans les
manuels d'introduction à la logique. En philosophie de la logique,
on a parfois considéré que les règles de la déduction naturelle donnent
le sens des opérateurs logiques : c'est la « sémantique fondée sur la
théorie de la démonstration ».

\end{document}
